\documentclass[aspectratio=169,11pt]{beamer}
% Lecture 04 — Collections, loops, ranges, do-while/split
% Course: Introduction to Programming with Kotlin

% Shared Beamer preamble for Introduction to Programming with Kotlin
% Include with: % Shared Beamer preamble for Introduction to Programming with Kotlin
% Include with: % Shared Beamer preamble for Introduction to Programming with Kotlin
% Include with: \input{preamble}

\usetheme{Madrid}
\usecolortheme{default}

% Clean, professional palette (deep navy + muted teal accent)
\definecolor{LectureNavy}{HTML}{1B3A4B}
\definecolor{LectureAccent}{HTML}{2A9D8F}
\definecolor{LectureGray}{HTML}{4A5568}
\definecolor{CodeBg}{HTML}{F7FAFC}
\definecolor{AlertBg}{HTML}{FFF5F5}
\definecolor{TryBg}{HTML}{F0FFF4}
\definecolor{QuizBg}{HTML}{EBF8FF}
\definecolor{AnswerKeyBg}{HTML}{FFFFF0}
\definecolor{AnswerGold}{HTML}{B7791F}
\definecolor{KeywordBlue}{HTML}{2B6CB0}
\definecolor{StringGreen}{HTML}{276749}
\definecolor{CommentGray}{HTML}{718096}

\setbeamercolor{structure}{fg=LectureNavy}
\setbeamercolor{palette primary}{bg=LectureNavy,fg=white}
\setbeamercolor{palette secondary}{bg=LectureAccent,fg=white}
\setbeamercolor{palette tertiary}{bg=LectureGray,fg=white}
% Madrid puts title/subtitle in a coloured box (inherits palette primary bg).
% Title text must be light on that dark box — not navy-on-navy.
\setbeamercolor{title}{fg=white,bg=LectureNavy}
\setbeamercolor{subtitle}{fg=white,bg=LectureNavy}
\setbeamercolor{author}{fg=LectureNavy}
\setbeamercolor{date}{fg=LectureGray}
\setbeamercolor{institute}{fg=LectureGray}
\setbeamercolor{frametitle}{bg=LectureNavy,fg=white}
\setbeamercolor{block title}{bg=LectureNavy,fg=white}
\setbeamercolor{block body}{bg=CodeBg,fg=black}
\setbeamercolor{block title alerted}{bg=LectureAccent,fg=white}
\setbeamercolor{block body alerted}{bg=TryBg,fg=black}
\setbeamercolor{block title example}{bg=LectureGray,fg=white}
\setbeamercolor{block body example}{bg=AlertBg,fg=black}

\setbeamertemplate{navigation symbols}{}
\setbeamertemplate{itemize items}[circle]
\setbeamertemplate{enumerate items}[default]

\usepackage[T1]{fontenc}
\usepackage[utf8]{inputenc}
\usepackage{lmodern}
\usepackage{booktabs}
\usepackage{tabularx}
\usepackage{graphicx}
\usepackage{hyperref}
\usepackage{listings}
\usepackage{xcolor}
\usepackage{tikz}
\usetikzlibrary{positioning,arrows.meta}

% listings: Kotlin without shell-escape (no built-in Kotlin language)
\lstdefinelanguage{Kotlin}{
  morekeywords={
    abstract,actual,annotation,as,break,by,catch,class,companion,const,
    constructor,continue,crossinline,data,do,dynamic,else,enum,expect,
    external,false,field,file,final,finally,for,fun,get,if,import,in,
    infix,init,inline,inner,interface,internal,is,lateinit,noinline,null,
    object,open,operator,out,override,package,private,protected,public,
    reified,return,sealed,set,super,suspend,tailrec,this,throw,true,try,
    typealias,typeof,val,var,when,where,while
  },
  sensitive=true,
  morecomment=[l]{//},
  morecomment=[s]{/*}{*/},
  morestring=[b]",
  morestring=[b]',
  morestring=[s]{"""}{"""},
}

\lstdefinestyle{kotlinstyle}{
  language=Kotlin,
  basicstyle=\ttfamily\small,
  keywordstyle=\color{KeywordBlue}\bfseries,
  commentstyle=\color{CommentGray}\itshape,
  stringstyle=\color{StringGreen},
  numbers=left,
  numberstyle=\tiny\color{CommentGray},
  numbersep=6pt,
  frame=single,
  rulecolor=\color{LectureGray!40},
  backgroundcolor=\color{CodeBg},
  breaklines=true,
  showstringspaces=false,
  tabsize=2,
  captionpos=b,
  morekeywords={readln,readLine,println,print,listOf,mutableListOf,toInt,toDouble,toLong,toShort,toByte,toChar,toString,split,size},
}
\lstset{style=kotlinstyle}

\newcommand{\kotlincode}[1]{\texttt{\small #1}}
\newcommand{\trythis}[1]{%
  \begin{alertblock}{Try this}
    #1
  \end{alertblock}%
}
\newcommand{\commonmistake}[1]{%
  \begin{exampleblock}{Common mistake}
    #1
  \end{exampleblock}%
}
% Formative in-class checks: question slide, then a dedicated Answer slide
\newcommand{\quizpause}{%
  \par\vspace{0.5em}
  {\small\textit{Pause --- answer (clicker / raise hand / neighbour) before the next slide.}}%
}
\newcommand{\quizbox}[1]{%
  \setbeamercolor{block title}{bg=KeywordBlue,fg=white}%
  \setbeamercolor{block body}{bg=QuizBg,fg=black}%
  \begin{block}{Quiz}
    #1
  \end{block}%
  \setbeamercolor{block title}{bg=LectureNavy,fg=white}%
  \setbeamercolor{block body}{bg=CodeBg,fg=black}%
}
\newcommand{\answerbox}[1]{%
  \setbeamercolor{block title}{bg=AnswerGold,fg=white}%
  \setbeamercolor{block body}{bg=AnswerKeyBg,fg=black}%
  \begin{block}{Answer}
    #1
  \end{block}%
  \setbeamercolor{block title}{bg=LectureNavy,fg=white}%
  \setbeamercolor{block body}{bg=CodeBg,fg=black}%
}

% Empty author/date placeholders for the lecturer to fill
\author{}
\date{}


\usetheme{Madrid}
\usecolortheme{default}

% Clean, professional palette (deep navy + muted teal accent)
\definecolor{LectureNavy}{HTML}{1B3A4B}
\definecolor{LectureAccent}{HTML}{2A9D8F}
\definecolor{LectureGray}{HTML}{4A5568}
\definecolor{CodeBg}{HTML}{F7FAFC}
\definecolor{AlertBg}{HTML}{FFF5F5}
\definecolor{TryBg}{HTML}{F0FFF4}
\definecolor{QuizBg}{HTML}{EBF8FF}
\definecolor{AnswerKeyBg}{HTML}{FFFFF0}
\definecolor{AnswerGold}{HTML}{B7791F}
\definecolor{KeywordBlue}{HTML}{2B6CB0}
\definecolor{StringGreen}{HTML}{276749}
\definecolor{CommentGray}{HTML}{718096}

\setbeamercolor{structure}{fg=LectureNavy}
\setbeamercolor{palette primary}{bg=LectureNavy,fg=white}
\setbeamercolor{palette secondary}{bg=LectureAccent,fg=white}
\setbeamercolor{palette tertiary}{bg=LectureGray,fg=white}
% Madrid puts title/subtitle in a coloured box (inherits palette primary bg).
% Title text must be light on that dark box — not navy-on-navy.
\setbeamercolor{title}{fg=white,bg=LectureNavy}
\setbeamercolor{subtitle}{fg=white,bg=LectureNavy}
\setbeamercolor{author}{fg=LectureNavy}
\setbeamercolor{date}{fg=LectureGray}
\setbeamercolor{institute}{fg=LectureGray}
\setbeamercolor{frametitle}{bg=LectureNavy,fg=white}
\setbeamercolor{block title}{bg=LectureNavy,fg=white}
\setbeamercolor{block body}{bg=CodeBg,fg=black}
\setbeamercolor{block title alerted}{bg=LectureAccent,fg=white}
\setbeamercolor{block body alerted}{bg=TryBg,fg=black}
\setbeamercolor{block title example}{bg=LectureGray,fg=white}
\setbeamercolor{block body example}{bg=AlertBg,fg=black}

\setbeamertemplate{navigation symbols}{}
\setbeamertemplate{itemize items}[circle]
\setbeamertemplate{enumerate items}[default]

\usepackage[T1]{fontenc}
\usepackage[utf8]{inputenc}
\usepackage{lmodern}
\usepackage{booktabs}
\usepackage{tabularx}
\usepackage{graphicx}
\usepackage{hyperref}
\usepackage{listings}
\usepackage{xcolor}
\usepackage{tikz}
\usetikzlibrary{positioning,arrows.meta}

% listings: Kotlin without shell-escape (no built-in Kotlin language)
\lstdefinelanguage{Kotlin}{
  morekeywords={
    abstract,actual,annotation,as,break,by,catch,class,companion,const,
    constructor,continue,crossinline,data,do,dynamic,else,enum,expect,
    external,false,field,file,final,finally,for,fun,get,if,import,in,
    infix,init,inline,inner,interface,internal,is,lateinit,noinline,null,
    object,open,operator,out,override,package,private,protected,public,
    reified,return,sealed,set,super,suspend,tailrec,this,throw,true,try,
    typealias,typeof,val,var,when,where,while
  },
  sensitive=true,
  morecomment=[l]{//},
  morecomment=[s]{/*}{*/},
  morestring=[b]",
  morestring=[b]',
  morestring=[s]{"""}{"""},
}

\lstdefinestyle{kotlinstyle}{
  language=Kotlin,
  basicstyle=\ttfamily\small,
  keywordstyle=\color{KeywordBlue}\bfseries,
  commentstyle=\color{CommentGray}\itshape,
  stringstyle=\color{StringGreen},
  numbers=left,
  numberstyle=\tiny\color{CommentGray},
  numbersep=6pt,
  frame=single,
  rulecolor=\color{LectureGray!40},
  backgroundcolor=\color{CodeBg},
  breaklines=true,
  showstringspaces=false,
  tabsize=2,
  captionpos=b,
  morekeywords={readln,readLine,println,print,listOf,mutableListOf,toInt,toDouble,toLong,toShort,toByte,toChar,toString,split,size},
}
\lstset{style=kotlinstyle}

\newcommand{\kotlincode}[1]{\texttt{\small #1}}
\newcommand{\trythis}[1]{%
  \begin{alertblock}{Try this}
    #1
  \end{alertblock}%
}
\newcommand{\commonmistake}[1]{%
  \begin{exampleblock}{Common mistake}
    #1
  \end{exampleblock}%
}
% Formative in-class checks: question slide, then a dedicated Answer slide
\newcommand{\quizpause}{%
  \par\vspace{0.5em}
  {\small\textit{Pause --- answer (clicker / raise hand / neighbour) before the next slide.}}%
}
\newcommand{\quizbox}[1]{%
  \setbeamercolor{block title}{bg=KeywordBlue,fg=white}%
  \setbeamercolor{block body}{bg=QuizBg,fg=black}%
  \begin{block}{Quiz}
    #1
  \end{block}%
  \setbeamercolor{block title}{bg=LectureNavy,fg=white}%
  \setbeamercolor{block body}{bg=CodeBg,fg=black}%
}
\newcommand{\answerbox}[1]{%
  \setbeamercolor{block title}{bg=AnswerGold,fg=white}%
  \setbeamercolor{block body}{bg=AnswerKeyBg,fg=black}%
  \begin{block}{Answer}
    #1
  \end{block}%
  \setbeamercolor{block title}{bg=LectureNavy,fg=white}%
  \setbeamercolor{block body}{bg=CodeBg,fg=black}%
}

% Empty author/date placeholders for the lecturer to fill
\author{}
\date{}


\usetheme{Madrid}
\usecolortheme{default}

% Clean, professional palette (deep navy + muted teal accent)
\definecolor{LectureNavy}{HTML}{1B3A4B}
\definecolor{LectureAccent}{HTML}{2A9D8F}
\definecolor{LectureGray}{HTML}{4A5568}
\definecolor{CodeBg}{HTML}{F7FAFC}
\definecolor{AlertBg}{HTML}{FFF5F5}
\definecolor{TryBg}{HTML}{F0FFF4}
\definecolor{QuizBg}{HTML}{EBF8FF}
\definecolor{AnswerKeyBg}{HTML}{FFFFF0}
\definecolor{AnswerGold}{HTML}{B7791F}
\definecolor{KeywordBlue}{HTML}{2B6CB0}
\definecolor{StringGreen}{HTML}{276749}
\definecolor{CommentGray}{HTML}{718096}

\setbeamercolor{structure}{fg=LectureNavy}
\setbeamercolor{palette primary}{bg=LectureNavy,fg=white}
\setbeamercolor{palette secondary}{bg=LectureAccent,fg=white}
\setbeamercolor{palette tertiary}{bg=LectureGray,fg=white}
% Madrid puts title/subtitle in a coloured box (inherits palette primary bg).
% Title text must be light on that dark box — not navy-on-navy.
\setbeamercolor{title}{fg=white,bg=LectureNavy}
\setbeamercolor{subtitle}{fg=white,bg=LectureNavy}
\setbeamercolor{author}{fg=LectureNavy}
\setbeamercolor{date}{fg=LectureGray}
\setbeamercolor{institute}{fg=LectureGray}
\setbeamercolor{frametitle}{bg=LectureNavy,fg=white}
\setbeamercolor{block title}{bg=LectureNavy,fg=white}
\setbeamercolor{block body}{bg=CodeBg,fg=black}
\setbeamercolor{block title alerted}{bg=LectureAccent,fg=white}
\setbeamercolor{block body alerted}{bg=TryBg,fg=black}
\setbeamercolor{block title example}{bg=LectureGray,fg=white}
\setbeamercolor{block body example}{bg=AlertBg,fg=black}

\setbeamertemplate{navigation symbols}{}
\setbeamertemplate{itemize items}[circle]
\setbeamertemplate{enumerate items}[default]

\usepackage[T1]{fontenc}
\usepackage[utf8]{inputenc}
\usepackage{lmodern}
\usepackage{booktabs}
\usepackage{tabularx}
\usepackage{graphicx}
\usepackage{hyperref}
\usepackage{listings}
\usepackage{xcolor}
\usepackage{tikz}
\usetikzlibrary{positioning,arrows.meta}

% listings: Kotlin without shell-escape (no built-in Kotlin language)
\lstdefinelanguage{Kotlin}{
  morekeywords={
    abstract,actual,annotation,as,break,by,catch,class,companion,const,
    constructor,continue,crossinline,data,do,dynamic,else,enum,expect,
    external,false,field,file,final,finally,for,fun,get,if,import,in,
    infix,init,inline,inner,interface,internal,is,lateinit,noinline,null,
    object,open,operator,out,override,package,private,protected,public,
    reified,return,sealed,set,super,suspend,tailrec,this,throw,true,try,
    typealias,typeof,val,var,when,where,while
  },
  sensitive=true,
  morecomment=[l]{//},
  morecomment=[s]{/*}{*/},
  morestring=[b]",
  morestring=[b]',
  morestring=[s]{"""}{"""},
}

\lstdefinestyle{kotlinstyle}{
  language=Kotlin,
  basicstyle=\ttfamily\small,
  keywordstyle=\color{KeywordBlue}\bfseries,
  commentstyle=\color{CommentGray}\itshape,
  stringstyle=\color{StringGreen},
  numbers=left,
  numberstyle=\tiny\color{CommentGray},
  numbersep=6pt,
  frame=single,
  rulecolor=\color{LectureGray!40},
  backgroundcolor=\color{CodeBg},
  breaklines=true,
  showstringspaces=false,
  tabsize=2,
  captionpos=b,
  morekeywords={readln,readLine,println,print,listOf,mutableListOf,toInt,toDouble,toLong,toShort,toByte,toChar,toString,split,size},
}
\lstset{style=kotlinstyle}

\newcommand{\kotlincode}[1]{\texttt{\small #1}}
\newcommand{\trythis}[1]{%
  \begin{alertblock}{Try this}
    #1
  \end{alertblock}%
}
\newcommand{\commonmistake}[1]{%
  \begin{exampleblock}{Common mistake}
    #1
  \end{exampleblock}%
}
% Formative in-class checks: question slide, then a dedicated Answer slide
\newcommand{\quizpause}{%
  \par\vspace{0.5em}
  {\small\textit{Pause --- answer (clicker / raise hand / neighbour) before the next slide.}}%
}
\newcommand{\quizbox}[1]{%
  \setbeamercolor{block title}{bg=KeywordBlue,fg=white}%
  \setbeamercolor{block body}{bg=QuizBg,fg=black}%
  \begin{block}{Quiz}
    #1
  \end{block}%
  \setbeamercolor{block title}{bg=LectureNavy,fg=white}%
  \setbeamercolor{block body}{bg=CodeBg,fg=black}%
}
\newcommand{\answerbox}[1]{%
  \setbeamercolor{block title}{bg=AnswerGold,fg=white}%
  \setbeamercolor{block body}{bg=AnswerKeyBg,fg=black}%
  \begin{block}{Answer}
    #1
  \end{block}%
  \setbeamercolor{block title}{bg=LectureNavy,fg=white}%
  \setbeamercolor{block body}{bg=CodeBg,fg=black}%
}

% Empty author/date placeholders for the lecturer to fill
\author{}
\date{}


\title{Lecture 04: Collections and Loops}
\subtitle{Introduction to Programming with Kotlin}

\begin{document}

%----------------------------------------------------------------------
    \begin{frame}
        \titlepage
    \end{frame}

%----------------------------------------------------------------------
    \begin{frame}{Agenda}
        \begin{enumerate}
            \item Introduction to collections
            \item \kotlincode{repeat} and \kotlincode{while}
            \item \kotlincode{for} and the \kotlincode{size} property
            \item Ranges
            \item \kotlincode{do-while} and \kotlincode{split}
        \end{enumerate}
    \end{frame}

%======================================================================
    \section{Introduction to collections}
%======================================================================

    \begin{frame}{Why collections?}
        So far: one value per variable.
        Often you need \textbf{many} values of the same kind.
        \vspace{0.6em}
        \begin{itemize}
            \item A \textbf{list} stores an ordered sequence of elements
            \item Access by index: first element at index \kotlincode{0}
            \item Today: read-only lists with \kotlincode{listOf}, and a peek at
              \kotlincode{mutableListOf}
        \end{itemize}
    \end{frame}

    \begin{frame}[fragile]{Creating a list}
        \begin{lstlisting}
fun main() {
    val nums = listOf(10, 20, 30)
    val names = listOf("Ada", "Grace", "Alan")

    println(nums[0])      // 10
    println(names[2])     // Alan
    println(nums.size)    // 3
}
        \end{lstlisting}
        \vspace{0.3em}
        \kotlincode{listOf(...)} creates a read-only list.
        Indexing uses \kotlincode{[i]}; length is \kotlincode{.size}.
    \end{frame}

    \begin{frame}[fragile]{Mutable lists (brief)}
        \begin{lstlisting}
fun main() {
    val scores = mutableListOf(5, 7)
    scores.add(9)
    scores[0] = 6
    println(scores)       // [6, 7, 9]
}
        \end{lstlisting}
        \vspace{0.3em}
        \begin{itemize}
            \item \kotlincode{mutableListOf} --- you can add/change elements
            \item The variable can still be \kotlincode{val} (the list object changes)
            \item Prefer read-only lists until mutation is needed
        \end{itemize}
    \end{frame}

    \begin{frame}{Indexing rules}
        \begin{itemize}
            \item Valid indices: \kotlincode{0} \ldots\ \kotlincode{size - 1}
            \item \kotlincode{list[list.size]} --- out of bounds (crash)
            \item Empty list: \kotlincode{size == 0}; no valid index
        \end{itemize}
        \vspace{0.5em}
        \commonmistake{Thinking the last index is \kotlincode{size}.
        It is \kotlincode{size - 1}.}
    \end{frame}

    %--- Quiz checkpoint 1 ---
    \begin{frame}{Quiz 1 --- Lists}
        \quizbox{For \kotlincode{val xs = listOf("a", "b", "c")}, what is \kotlincode{xs[1]}?}
        \begin{enumerate}
            \item \kotlincode{"a"}
            \item \kotlincode{"b"}
            \item \kotlincode{"c"}
            \item \kotlincode{1}
        \end{enumerate}
        \quizpause
    \end{frame}

    \begin{frame}{Answer --- Quiz 1}
        \answerbox{\textbf{2.} \kotlincode{"b"}}
        \vspace{0.4em}
        Indexing is zero-based: \kotlincode{0 → "a"}, \kotlincode{1 → "b"}, \kotlincode{2 → "c"}.
    \end{frame}

%======================================================================
    \section{repeat and while}
%======================================================================

    \begin{frame}{Repeating work}
        Loops run code multiple times.
        \vspace{0.5em}
        \begin{itemize}
            \item \kotlincode{repeat(n) \{ \ldots \}} --- do something exactly $n$ times
            \item \kotlincode{while (condition) \{ \ldots \}} --- while the condition stays true
        \end{itemize}
    \end{frame}

    \begin{frame}[fragile]{\kotlincode{repeat}}
        \begin{lstlisting}
fun main() {
    repeat(3) {
        println("Hello")
    }
    repeat(5) { i ->
        println("Index $i")  // 0, 1, 2, 3, 4
    }
}
        \end{lstlisting}
        \vspace{0.3em}
        The optional parameter is the iteration index starting at $0$.
    \end{frame}

    \begin{frame}[fragile]{\kotlincode{while}}
        \begin{lstlisting}
fun main() {
    var n = readln().toInt()
    while (n > 0) {
        println(n)
        n -= 1
    }
    println("Done")
}
        \end{lstlisting}
        \vspace{0.3em}
        Check the condition \textbf{before} each iteration.
        If it starts false, the body never runs.
    \end{frame}

    \begin{frame}[fragile]{Avoid infinite loops}
        \begin{lstlisting}
// Dangerous: n never changes
// while (n > 0) {
//     println(n)
// }
        \end{lstlisting}
        \vspace{0.3em}
        Ensure something in the loop moves toward making the condition false
        (counter update, new input, \kotlincode{break}, \ldots).
        \vspace{0.4em}
        \trythis{Read integers until the user enters \kotlincode{0}, summing the rest.}
    \end{frame}

    %--- Quiz checkpoint 2 ---
    \begin{frame}{Quiz 2 --- while}
        \quizbox{How many times does the body run?\\
        \kotlincode{var i = 0; while (i < 3) \{ i += 1 \}}}
        \begin{enumerate}
            \item 0
            \item 2
            \item 3
            \item Forever
        \end{enumerate}
        \quizpause
    \end{frame}

    \begin{frame}{Answer --- Quiz 2}
        \answerbox{\textbf{3.} 3 times}
        \vspace{0.4em}
        Iterations with \kotlincode{i} equal to $0$, $1$, then $2$.
        After the third increment, \kotlincode{i == 3} and the loop stops.
    \end{frame}

%======================================================================
    \section{for and size}
%======================================================================

    \begin{frame}{\kotlincode{for} over a collection}
        The most common loop over lists:
        \vspace{0.4em}
        \begin{itemize}
            \item \kotlincode{for (item in list) \{ \ldots \}}
            \item Visit each element in order
            \item No manual index required
        \end{itemize}
    \end{frame}

    \begin{frame}[fragile]{Iterating elements}
        \begin{lstlisting}
fun main() {
    val names = listOf("Ada", "Grace", "Alan")
    for (name in names) {
        println("Hello, $name")
    }
}
        \end{lstlisting}
    \end{frame}

    \begin{frame}[fragile]{Using \kotlincode{size} with indices}
        \begin{lstlisting}
fun main() {
    val nums = listOf(4, 8, 15, 16)
    println("Count: ${nums.size}")

    for (i in 0 until nums.size) {
        println("nums[$i] = ${nums[i]}")
    }
    // Or: for (i in nums.indices) { ... }
}
        \end{lstlisting}
        \vspace{0.3em}
        \kotlincode{0 until size} is $0, 1, \ldots, size-1$ --- safe indices.
    \end{frame}

    \begin{frame}[fragile]{Summing a list}
        \begin{lstlisting}
fun main() {
    val xs = listOf(1, 2, 3, 4)
    var sum = 0
    for (x in xs) {
        sum += x
    }
    println(sum)  // 10
}
        \end{lstlisting}
        \trythis{Find the maximum element of a list of integers using a \kotlincode{for} loop.}
    \end{frame}

    %--- Quiz checkpoint 3 ---
    \begin{frame}{Quiz 3 --- for and size}
        \quizbox{For a list with \kotlincode{size == 4}, which index range is safe?}
        \begin{enumerate}
            \item \kotlincode{1..4}
            \item \kotlincode{0..4}
            \item \kotlincode{0 until 4}
            \item \kotlincode{0..size}
        \end{enumerate}
        \quizpause
    \end{frame}

    \begin{frame}{Answer --- Quiz 3}
        \answerbox{\textbf{3.} \kotlincode{0 until 4}}
        \vspace{0.4em}
        Valid indices are $0,1,2,3$.
        \kotlincode{0 until 4} produces exactly those values;
        \kotlincode{0..4} would include $4$ (out of bounds).
    \end{frame}

%======================================================================
    \section{Ranges}
%======================================================================

    \begin{frame}{Ranges}
        A \textbf{range} is a sequence of values from a start to an end.
        \vspace{0.5em}
        \begin{itemize}
            \item \kotlincode{1..5} --- inclusive: $1,2,3,4,5$
            \item \kotlincode{1 until 5} --- exclusive end: $1,2,3,4$
            \item \kotlincode{5 downTo 1} --- descending
            \item \kotlincode{1..10 step 2} --- $1,3,5,7,9$
        \end{itemize}
    \end{frame}

    \begin{frame}[fragile]{Looping over ranges}
        \begin{lstlisting}
fun main() {
    for (i in 1..5) {
        print("$i ")
    }
    println()
    for (i in 5 downTo 1) {
        print("$i ")
    }
    println()
    for (i in 0 until 3) {
        print("$i ")
    }
}
        \end{lstlisting}
        Output idea: \kotlincode{1 2 3 4 5} then \kotlincode{5 4 3 2 1} then \kotlincode{0 1 2}.
    \end{frame}

    \begin{frame}[fragile]{\kotlincode{in} for membership}
        \begin{lstlisting}
fun main() {
    val age = readln().toInt()
    if (age in 13..19) {
        println("Teen")
    }
    val letter = readln()[0]
    if (letter in 'a'..'z') {
        println("Lowercase letter")
    }
}
        \end{lstlisting}
        Ranges work for integers and characters.
    \end{frame}

    %--- Quiz checkpoint 4 ---
    \begin{frame}{Quiz 4 --- Ranges}
        \quizbox{What values does \kotlincode{1 until 4} contain?}
        \begin{enumerate}
            \item $1,2,3,4$
            \item $1,2,3$
            \item $0,1,2,3$
            \item $1,4$
        \end{enumerate}
        \quizpause
    \end{frame}

    \begin{frame}{Answer --- Quiz 4}
        \answerbox{\textbf{2.} $1,2,3$}
        \vspace{0.4em}
        \kotlincode{until} excludes the end value. Inclusive form is \kotlincode{1..4}.
    \end{frame}

%======================================================================
    \section{do-while and split}
%======================================================================

    \begin{frame}{\kotlincode{do-while}}
        Like \kotlincode{while}, but the body runs \textbf{at least once},
        then the condition is checked.
        \vspace{0.5em}
        Useful for ``ask, then decide whether to ask again'' menus and input loops.
    \end{frame}

    \begin{frame}[fragile]{\kotlincode{do-while} example}
        \begin{lstlisting}
fun main() {
    var answer: String
    do {
        print("Type 'quit' to stop: ")
        answer = readln()
        println("You typed: $answer")
    } while (answer != "quit")
}
        \end{lstlisting}
        \vspace{0.3em}
        Compare with \kotlincode{while}: here the prompt always appears at least once.
    \end{frame}

    \begin{frame}{\kotlincode{split} on strings}
        \kotlincode{String.split(...)} breaks text into a list of pieces.
        \vspace{0.5em}
        \begin{itemize}
            \item Common: split on spaces or commas
            \item Result type: \kotlincode{List<String>}
            \item Great for reading several values from one line
        \end{itemize}
    \end{frame}

    \begin{frame}[fragile]{Reading several numbers on one line}
        \begin{lstlisting}
fun main() {
    print("Enter numbers separated by spaces: ")
    val parts = readln().split(" ")
    var sum = 0
    for (p in parts) {
        if (p.isNotEmpty()) {
            sum += p.toInt()
        }
    }
    println("Sum = $sum")
}
        \end{lstlisting}
        \commonmistake{Forgetting that \kotlincode{split} returns strings --- still need
        \kotlincode{toInt()} for arithmetic.}
    \end{frame}

    \begin{frame}[fragile]{Split on commas}
        \begin{lstlisting}
fun main() {
    val line = "Ada,Grace,Alan"
    val names = line.split(",")
    for (name in names) {
        println(name.trim())
    }
}
        \end{lstlisting}
        \kotlincode{trim()} removes surrounding spaces after splitting.
        \vspace{0.4em}
        \trythis{Read \kotlincode{a,b,c} as three integers from one line and print the average.}
    \end{frame}

    %--- Quiz checkpoint 5 ---
    \begin{frame}{Quiz 5 --- split}
        \quizbox{What is \kotlincode{"10 20 30".split(" ").size}?}
        \begin{enumerate}
            \item 1
            \item 2
            \item 3
            \item 10
        \end{enumerate}
        \quizpause
    \end{frame}

    \begin{frame}{Answer --- Quiz 5}
        \answerbox{\textbf{3.} 3}
        \vspace{0.4em}
        Splitting on spaces yields three strings: \kotlincode{"10"}, \kotlincode{"20"},
        \kotlincode{"30"}.
    \end{frame}

%======================================================================
    \section{Wrap-up}
%======================================================================

    \begin{frame}{Summary}
        \begin{itemize}
            \item Lists: \kotlincode{listOf}, index from $0$, property \kotlincode{size}
            \item \kotlincode{repeat} / \kotlincode{while} / \kotlincode{for} / \kotlincode{do-while}
            \item Ranges: \kotlincode{..}, \kotlincode{until}, \kotlincode{downTo}, \kotlincode{step}
            \item \kotlincode{split} turns a string into a list of pieces
        \end{itemize}
    \end{frame}

    \begin{frame}{Looking ahead}
        Next: the \kotlincode{Char} type --- single characters, codes, and how they
        relate to strings.
    \end{frame}

\end{document}
